\documentclass[11pt]{article}

\usepackage[margin=1in]{geometry}
\usepackage{amsmath,amssymb,amsthm}
\usepackage{bm}
\usepackage{mathtools}

\newtheorem{lemma}{Lemma}
\newtheorem{proposition}{Proposition}
\newtheorem{remark}{Remark}

\newcommand{\bbeta}{\bm{\beta}}
\newcommand{\bhat}{\widehat{\beta}}
\newcommand{\pstar}{p^{\star}}
\newcommand{\lstar}{\lambda^{\star}}
\newcommand{\thetahat}{\widehat{\theta}}
\newcommand{\Estar}{\mathrm{E}}
\newcommand{\lamdiff}{(\lambda_0-\lambda_1)}
\newcommand{\betasub}{\bhat_{\backslash j}}
\DeclareMathOperator*{\argmax}{arg\,max}

\title{The Zero Gap of the Spike-and-Slab LASSO Global Mode:\\
An Expanded Derivation and its Relation to the Adaptive Threshold}
\author{Sophie Huebler}
\date{\today}

\begin{document}
\maketitle

\noindent
This note expands the remark at the end of Section~3.1 of Ro\v{c}kov\'a and
George (2018), \emph{The Spike-and-Slab LASSO}. We prove that the nonzero
coordinates of the global mode are bounded away from zero (the \emph{zero
gap}), and we relate this magnitude floor $\delta_j$ to the adaptive selection
threshold $\Delta_j$ of Theorem~4. Equation numbers in parentheses refer to
that paper.

\section{The claim}

Under the conditions of Theorem~4, the global mode $\bhat$ has a \emph{zero
gap}: every coordinate is either exactly zero or bounded away from zero, with
\begin{equation}\label{eq:claim}
  |\bhat_j| > \delta_j
  \qquad\text{whenever } \bhat_j \neq 0,
\end{equation}
where $\delta_j$ is the unique solution of $\pstar_{\thetahat_j}(\delta_j)=c_+$,
with $\thetahat_j = \Estar[\theta \mid \betasub]$ and
\begin{equation}\label{eq:cplus}
  c_+ = \tfrac{1}{2}\Bigl(1 + \sqrt{1 - 4n/\lamdiff^2}\,\Bigr).
\end{equation}
This is the nonseparable analogue of the Section~2 statement
$|\bhat_j| > \delta_{c_+}$ (attributed to Lemma~4.1 of Ro\v{c}kov\'a 2015).
Because the coordinatewise problem freezes $\theta$ at $\thetahat_j$, the proof
is identical to the separable one with $\theta = \thetahat_j$; the only new
feature is that the threshold is coordinate-specific and data-driven.

\section{Univariate reduction}

By (3.12), with $z=z_j>0$ fixed (WLOG, so $\bhat_j\ge 0$) and
$\theta=\thetahat_j$ held fixed,
\begin{equation}\label{eq:uniobj}
  \bhat_j = \argmax_{\beta}\; Q(\beta),
  \qquad
  Q(\beta) = -\tfrac{1}{2}(z-n\beta)^2 + n\,\rho(\beta\mid\theta),
\end{equation}
where $\rho$ is the separable singleton penalty (2.2),
\begin{equation}\label{eq:rho}
  \rho(\beta\mid\theta)
    = -\lambda_1|\beta| + \log\!\bigl[\pstar_\theta(0)/\pstar_\theta(\beta)\bigr].
\end{equation}

\section{Derivatives and the curvature of the penalty}

From Lemma~1, $\rho'(\beta) = -\lstar_\theta(\beta)$ with
$\lstar_\theta(\beta) = \lambda_1\pstar_\theta(\beta) + \lambda_0[1-\pstar_\theta(\beta)]
= \lambda_0 - \lamdiff\,\pstar_\theta(\beta)$.
Differentiating the exponential mixing weight (2.6) gives the logistic identity
\begin{equation}\label{eq:pprime}
  \pstar_\theta{}'(\beta) = \lamdiff\,\pstar_\theta(\beta)\bigl[1-\pstar_\theta(\beta)\bigr],
\end{equation}
so the penalty's second derivative is
\begin{equation}\label{eq:rhosecond}
  \rho''(\beta) = \lamdiff^2\,\pstar_\theta(\beta)\bigl[1-\pstar_\theta(\beta)\bigr].
\end{equation}
Its maximum over $\beta$ is the maximal nonconcavity
$\kappa = \tfrac14\lamdiff^2$ (attained at $\pstar_\theta=\tfrac12$), matching
p.~433. Hence
\begin{equation}\label{eq:Qderiv}
  Q'(\beta) = n\bigl[z - n\beta - \lstar_\theta(\beta)\bigr],
  \qquad
  Q''(\beta) = n\bigl[\rho''(\beta) - n\bigr].
\end{equation}

\section{The second-order condition forces $\bhat$ onto the upper branch}

If $\bhat>0$ is the global maximizer it is in particular a local maximizer, so
$Q''(\bhat)\le 0$, i.e.\ $\rho''(\bhat)\le n$. With $u := \pstar_\theta(\bhat)$
and \eqref{eq:rhosecond},
\begin{equation}\label{eq:soc}
  \lamdiff^2\, u(1-u) \le n
  \iff
  u(1-u) \le \frac{n}{\lamdiff^2}
  \iff
  u^2 - u + \frac{n}{\lamdiff^2} \ge 0.
\end{equation}
The quadratic in \eqref{eq:soc} is nonnegative iff
\begin{equation}\label{eq:branches}
  u \le c_-
  \quad\text{or}\quad
  u \ge c_+,
  \qquad
  c_\mp = \tfrac{1}{2}\Bigl(1 \mp \sqrt{1-4n/\lamdiff^2}\Bigr),
\end{equation}
with real, distinct roots precisely when $\lamdiff^2 > 4n$. Thus no local
maximum can lie in the curvature band $c_- < \pstar_\theta(\beta) < c_+$, where
$Q$ is strictly convex ($Q''>0$). Consequently $Q$ is concave--convex--concave
in $\beta$, and the only nonzero maximizer candidates are a point on the
\emph{lower} branch ($\pstar_\theta \le c_-$, small $\beta$) or the \emph{upper}
branch ($\pstar_\theta \ge c_+$, larger $\beta$).

\paragraph{Excluding the lower branch.}
Write $h(\beta) := Q'(\beta)/n = z - n\beta - \lstar_\theta(\beta)$, so by
\eqref{eq:Qderiv} $h'(\beta) = \rho''(\beta) - n$ shares the sign pattern of
\eqref{eq:branches}: $h$ decreases on $(0,\delta_{c_-})$, increases across the
convex band, and decreases again for $\beta > \delta_{c_+}$. Near the origin the
spike slope dominates, $\lstar_\theta(0)\approx\lambda_0$, so $Q$ falls away
from $0$ across the first concave piece and any lower-branch stationary point is
dominated by the mode at $0$. The condition $g_\theta(0)>0$ --- equivalently
$\lstar_\theta(0) > \sqrt{2n\log[1/\pstar_\theta(0)]} + \lambda_1$ (p.~435) ---
is exactly what guarantees this domination, so the only nonzero value that can
beat $Q(0)$ lies on the upper decreasing branch. (This is the content of
Lemma~4.1 of Ro\v{c}kov\'a 2015; the second-order argument above is the crux.)
Hence a nonzero global mode satisfies $\pstar_\theta(\bhat)\ge c_+$.

\section{Converting the branch condition into the magnitude floor}

Since $\pstar_\theta$ is strictly increasing,
$\pstar_\theta(\bhat)\ge c_+ \iff \bhat \ge \delta_j$, where $\delta_j$ solves
$\pstar_{\thetahat_j}(\delta_j)=c_+$. Solving (2.6) explicitly,
\begin{equation*}
  c_+ = \Bigl[1 + \tfrac{\lambda_0}{\lambda_1}\tfrac{1-\theta}{\theta}
        e^{-\delta_j\lamdiff}\Bigr]^{-1},
\end{equation*}
gives
\begin{equation}\label{eq:delta}
  \boxed{\;
  \delta_j = \frac{1}{\lamdiff}\,
    \log\!\left[
      \frac{1-\thetahat_j}{\thetahat_j}\cdot
      \frac{\lambda_0}{\lambda_1}\cdot
      \frac{c_+}{1-c_+}
    \right].
  \;}
\end{equation}
With a single fixed $\theta$ this is exactly $\delta_{c_+}$ of Section~2; with
$\thetahat_j = \Estar[\theta\mid\betasub]$ it becomes the coordinate-specific
gap. Because $z>\Delta_j$ at selection places the stationary point strictly
interior to the upper branch, the inequality is strict, establishing
$|\bhat_j| > \delta_j$ and hence \eqref{eq:claim}. \qed

\section{Relation to the adaptive threshold $\Delta_j$ of Theorem~4}

The two thresholds answer different questions and arise from different parts of
the same geometry.

\medskip
\noindent\textbf{$\Delta_j$ is a threshold on the data $z_j$ --- it governs
\emph{selection}.} By Theorem~3, $\bhat_j = 0 \iff |z_j|\le\Delta_j$, and
Theorem~4 bounds it by $\Delta_j^L < \Delta_j < \Delta_j^U$ with
$\Delta_j^U = \sqrt{2n\log[1/\pstar_{\thetahat_j}(0)]} + \lambda_1$. It is a
\emph{global} comparison --- the value of $z_j$ at which $Q(\bhat)$ overtakes
$Q(0)$ --- which is why its scale is $\sqrt{2n\log[1/\pstar_{\thetahat_j}(0)]}$
and it is driven by $\pstar_{\thetahat_j}(0)$, the shrinkage at the origin.

\medskip
\noindent\textbf{$\delta_j$ is a threshold on the estimate $\bhat_j$ --- it
governs the \emph{magnitude} of whatever survives selection.} It is a
\emph{local} condition: the boundary $Q''=0$ of the concave region, which is why
its scale is set by $c_+$, i.e.\ by $\lamdiff^2$ versus $n$, rather than by
$\pstar_{\thetahat_j}(0)$.

\medskip
They are tied together in three ways.

\begin{enumerate}
\item \textbf{Same adaptive driver.} Both inherit coordinate-adaptivity from the
single learned quantity $\thetahat_j = \Estar[\theta\mid\betasub]$. As more
sparsity is detected in $\betasub$, $\thetahat_j$ falls; this raises
$\Delta_j^U$ (a smaller $\pstar_{\thetahat_j}(0)$, harder to enter the model)
while simultaneously raising $\delta_j$ (through the
$(1-\thetahat_j)/\thetahat_j$ factor in \eqref{eq:delta}). The
Scott--Berger multiplicity adjustment thus moves the selection threshold and the
gap \emph{in concert}, coordinate by coordinate --- the nonseparable counterpart
of the single fixed pair $(\Delta^U,\delta_{c_+})$ in Section~2.

\item \textbf{The selection threshold \emph{creates} the gap.} At the exact
selection boundary $z_j = \Delta_j$ the competing nonzero maximizer is already on
the upper branch ($\pstar_\theta \ge c_+$). So as $z_j$ crosses $\Delta_j$ the
estimate jumps discontinuously from $0$ to a value already exceeding $\delta_j$
--- the hard-thresholding face of the ``blend of soft- and hard-thresholding''
in Theorem~3. There is no continuum of tiny nonzero estimates; that is precisely
the zero gap.

\item \textbf{Common large-$\lambda_0$ limit.} As $\lamdiff\to\infty$,
$c_+\to 1$ and $1-c_+ = c_- \to n/\lamdiff^2$, so
$c_+/(1-c_+) \approx \lamdiff^2/n$ and
\begin{equation*}
  \delta_j \approx \frac{1}{\lamdiff}
    \log\!\left[
      \frac{1-\thetahat_j}{\thetahat_j}\cdot
      \frac{\lambda_0}{\lambda_1}\cdot
      \frac{\lamdiff^2}{n}
    \right] \longrightarrow 0.
\end{equation*}
The gap floor vanishes while $\Delta_j\to\Delta_j^U$, so the estimator
approaches the $\ell_0$/hard-threshold limit in which selection is controlled
entirely by $\Delta_j$ and surviving coefficients are essentially unbiased ---
matching the Section~2 remark that $\Delta\to\Delta^U$ for large $\lamdiff$.
\end{enumerate}

\medskip
\noindent
In short: $\Delta_j$ (from the \emph{global} value comparison, scale
$\sqrt{2n\log[1/\pstar_{\thetahat_j}(0)]}$) decides \emph{whether} a coordinate
is nonzero; $\delta_j$ (from the \emph{local} curvature boundary, scale set by
$c_+$) decides \emph{how small a nonzero coordinate can be}. Both are made
coordinate-adaptive by $\thetahat_j$, and they meet at the selection boundary,
which is what makes the gap appear.

\appendix

\section{The penalty derivative: a detailed derivation of Lemma~1}
\label{app:lemma1}

The proof of Lemma~1 opens from the identity
\begin{equation}\label{eq:lem1start}
  \frac{\partial\,\mathrm{pen}_S(\bbeta\mid\theta)}{\partial|\beta_j|}
  = \pstar_\theta(\beta_j)\,
      \frac{\partial\log\psi_1(\beta_j)}{\partial|\beta_j|}
  + \bigl[1-\pstar_\theta(\beta_j)\bigr]\,
      \frac{\partial\log\psi_0(\beta_j)}{\partial|\beta_j|}.
\end{equation}
It is worth noting that \eqref{eq:lem1start} is \emph{not} the derivative of the
regrouped singleton (2.2); it is the derivative of the original mixture form
(2.1), and (2.2) is an algebraically equivalent regrouping of the same object.
We record both routes.

\paragraph{Setup.}
The two Laplace components (1.3) are
\begin{equation*}
  \psi_1(\beta) = \tfrac{\lambda_1}{2}e^{-\lambda_1|\beta|},
  \qquad
  \psi_0(\beta) = \tfrac{\lambda_0}{2}e^{-\lambda_0|\beta|},
\end{equation*}
each depending on $\beta$ only through $t=|\beta|$, so that
\begin{equation}\label{eq:logscore}
  \frac{\partial\log\psi_1}{\partial|\beta|} = -\lambda_1,
  \qquad
  \frac{\partial\log\psi_0}{\partial|\beta|} = -\lambda_0.
\end{equation}
The per-coordinate term of the penalty, \emph{before} the rewrite into (2.2), is
the mixture form (2.1),
\begin{equation*}
  s(\beta_j) = \log\bigl[\theta\psi_1(\beta_j)+(1-\theta)\psi_0(\beta_j)\bigr]
             - \log\bigl[\theta\psi_1(0)+(1-\theta)\psi_0(0)\bigr],
\end{equation*}
whose second bracket is constant in $\beta_j$.

\paragraph{Where \eqref{eq:lem1start} comes from.}
Write the mixture density $f(\beta) := \theta\psi_1(\beta)+(1-\theta)\psi_0(\beta)$.
Then, with ${}'$ denoting $\partial/\partial|\beta|$,
\begin{equation*}
  \frac{\partial s}{\partial|\beta_j|}
  = \frac{\partial\log f}{\partial|\beta_j|}
  = \frac{f'}{f}
  = \frac{\theta\psi_1'(\beta_j)+(1-\theta)\psi_0'(\beta_j)}{f}.
\end{equation*}
Applying the standard log-mixture derivative identity --- multiply and divide
each term by its own component density ---
\begin{equation*}
  \frac{f'}{f}
  = \frac{\theta\psi_1}{f}\cdot\frac{\psi_1'}{\psi_1}
  + \frac{(1-\theta)\psi_0}{f}\cdot\frac{\psi_0'}{\psi_0}
  = \frac{\theta\psi_1}{f}\,\frac{\partial\log\psi_1}{\partial|\beta_j|}
  + \frac{(1-\theta)\psi_0}{f}\,\frac{\partial\log\psi_0}{\partial|\beta_j|}.
\end{equation*}
The two weights are exactly the conditional mixing probability (2.4) and its
complement,
\begin{equation}\label{eq:responsibility}
  \frac{\theta\psi_1(\beta_j)}{f} = \pstar_\theta(\beta_j),
  \qquad
  \frac{(1-\theta)\psi_0(\beta_j)}{f} = 1-\pstar_\theta(\beta_j),
\end{equation}
which yields \eqref{eq:lem1start}. Thus the identity is the score of a
log-mixture, equal to the \emph{responsibility-weighted average of the component
scores}. Here $\pstar_\theta(\beta_j)$ is not incidental: it is the posterior
probability $\mathrm{P}(\gamma_j=1\mid\beta_j,\theta)$ of (2.6), i.e.\ the
responsibility that the slab $\psi_1$ generated $\beta_j$ --- the same weight
that appears in the E-step of EM --- which is why it drops out immediately.
Substituting \eqref{eq:logscore},
\begin{equation*}
  \frac{\partial\,\mathrm{pen}_S(\bbeta\mid\theta)}{\partial|\beta_j|}
  = -\lambda_1\pstar_\theta(\beta_j) - \lambda_0\bigl[1-\pstar_\theta(\beta_j)\bigr]
  = -\lstar_\theta(\beta_j),
\end{equation*}
which is (2.7). \qed

\paragraph{Reconciliation with the regrouped form (2.2).}
Differentiating (2.2) directly must give the same answer; it merely hides the
mixture step inside the $\pstar$ term. From
$\rho(\beta_j\mid\theta) = -\lambda_1|\beta_j| + \log\pstar_\theta(0)
- \log\pstar_\theta(\beta_j)$,
\begin{equation*}
  \frac{\partial\rho}{\partial|\beta_j|}
  = -\lambda_1 - \frac{\partial\log\pstar_\theta(\beta_j)}{\partial|\beta_j|}.
\end{equation*}
Evaluating the last derivative still requires the mixture: since
$\log\pstar_\theta = \log\theta + \log\psi_1 - \log f$,
\begin{equation*}
  \frac{\partial\log\pstar_\theta}{\partial|\beta_j|}
  = \frac{\partial\log\psi_1}{\partial|\beta_j|}
    - \frac{\partial\log f}{\partial|\beta_j|}
  = -\lambda_1 - \bigl(-\lstar_\theta(\beta_j)\bigr)
  = \lstar_\theta(\beta_j) - \lambda_1,
\end{equation*}
using $\partial\log f/\partial|\beta_j| = -\lstar_\theta(\beta_j)$ from above.
Hence
\begin{equation*}
  \frac{\partial\rho}{\partial|\beta_j|}
  = -\lambda_1 - \bigl(\lstar_\theta(\beta_j)-\lambda_1\bigr)
  = -\lstar_\theta(\beta_j),
\end{equation*}
in agreement. The two forms coincide because (2.2) is the ``LASSO piece plus
$\log\pstar$ correction'' regrouping of (2.1) (the passage from 2.1 to 2.3);
but differentiating (2.2) directly forces one to compute
$\partial\log\pstar_\theta/\partial|\beta_j|$, which reintroduces the mixture.
Differentiating the mixture form first lets the log-derivative identity
\eqref{eq:responsibility} collapse everything to the $\pstar$-weighted average
in a single line.

\end{document}
